\section{An existing Pell denominator removes the artificial scale}
\label{sec:compact-unscaled}

The compact branch gives two further fixed equations in ten strictly
positive integers. Their exact arithmetic cost and total degree are
\[
 (78,630)\qquad\text{and}\qquad(79,468).
\]
Here an operation is one live binary addition, subtraction, or
multiplication, including a square or multiplication by a fixed integer.
The complete arithmetic circuits are \code{compact/equation78.json} and
\code{compact/equation79.json}. These are cost--degree tradeoffs against
the pinned package 133's $(84,436)$ and $(85,328)$ equations. The number
of unknowns and the entire fixed coefficient tuple are unchanged.

The new idea is an order of deduction. A large known divisor of a
denominator first identifies a Pell index. The identified index then
restores divisibility by the whole denominator. Pairwise coprimality of
all carrier denominators is unnecessary. This allows us to delete the
artificial scale $\mu=q^2+1$ and reuse the existing square $Y^2$.
The quotient substitution $Y=q^3s$, Pell doubling, and the final
product-square equation have predecessors in the pinned repository.

\subsection{Literal polynomial and fixed coefficients}

The only unknowns are
\[
 J,F,\alpha,i,T,w,s,n,\eta,\zeta\in\Z_{>0}.
\]
The fixed integers $\beta,K_0,M_C,M_F$ are the complete directly indexed
BFS943 tuple of package 133. They satisfy
\[
 \beta\ge31\text{ odd},\qquad K_0>0,\qquad
 0<M_C<\beta<M_F<2\beta.
\]
The supplied compressed unsigned little-endian payloads fix every byte;
\code{compact/rebuild_coefficients.py} reconstructs all four integers
from the unchanged finite tile specification and its 436-character code.
All four reconstructed raw hashes agree with the pinned parent.
The largest fixed port has 592,884 bits. This is the size of the stored
constant, not the height of the fully expanded polynomial.

All following symbols are polynomial abbreviations:
\begin{gather*}
 q=\beta J+1,\quad h=q-1,\quad C=q-F-\alpha,\quad X=wq,\\
 Y=q^3s,\quad k=\eta+\zeta,\quad c=kY+\eta,\quad D=Y^2,\\
 A=Y(X+1)+1,\quad \Delta=A^2-1,\quad H=4A-5,\\
 V_0=XD,\qquad E=1+V_0(V_0+1)k^2,\\
 R=[h(q-F)+\alpha](q^2-1)+(M_C+qM_F)J,\\
 e=k-R-1,\quad B=(K_0+w)C-F,\quad U=ic^2,\\
 a=\Delta U^2E,\quad F_0=1+a,\quad G=EF_0,\\
 F_*=1+4aF_0,\qquad Q=\Delta(F_*-1),\\
 b=c(T+\alpha)-R,\quad V=F_*b-c,\quad K_+=X(3c-2X)+2.
\end{gather*}
Choose $L=V^2$ for the 78-operation equation, or $L=bV$ for the
79-operation equation, and define
\begin{align*}
 M&=hDH,\\
 Z&=hD\bigl[H(n-LB)+K_+\bigr]-H(DB+he),\\
 R_G&=M^2G.
\end{align*}
The single equation is
\begin{equation}
 \boxed{\ P_L=Q(Z^2-R_GV^2)+(R_G-Z^2)=0.\ }
 \label{eq:compact-unscaled}
\end{equation}
There are no additional quantified variables, nonpolynomial primitives,
or implicit conditions in this definition. The actual circuits compute
$A-1$, $\Delta=(A-1)((A-1)+2)$, and $H=4(A-1)-1$.

\begin{theorem}[Unscaled compact consistency equations]
For either stated choice of $L$, ordinary arithmetic formalizable in ZFC,
composed with the pinned indexed native/refutation calibration, proves
\[
 \ZFC\vdash
 \bigl(\neg\exists\mathbf x\in\Z_{>0}^{10}\;
 P_L(\mathbf x)=0\bigr)\ \longleftrightarrow\ \Con(\ZFC).
\]
The exact degrees are 630 and 468, and the supplied evaluators have 78
and 79 live binary operations, respectively.
\end{theorem}

\subsection{Inherited interfaces and their scope}

The proof uses the following explicitly retained ordinary arithmetic
interfaces. They are statements of earlier work, not claims made by a
finite coefficient checker.
\begin{enumerate}
\item The universal main/auxiliary Pell folding and rank theorems:
from the strong norm, $c^2\mid I$, and the fixed-minus auxiliary
residues, they first recover $c=\psi_A(p)$ for an odd positive $p$,
and, once $0<p,R<c/2$, recover $p=R$.
\item The universal projection and fixed-minus parity theorems:
after the required index and size bounds, they give the whole-product
projection modulo $4A-5$ and $R\equiv3\pmod4$.
\item The zero-offset ratio estimate, proved uniformly for every
$\lambda\ge1$, and its half-binomial rounding consequences.
\item The indexed native coefficient calibration and fresh source
reconstruction: the resulting outer arithmetic data are equivalent to
a finite syntactic ZFC refutation. This branch retains its native/K215
ancestry; it is separate from the direct-calculus quartics.
\end{enumerate}
The exact source files and immutable commit are recorded in
\code{compact/SOURCE_PINS.json}. The new argument below verifies the
changed premises and supplies both existential directions. It does not
claim a new end-to-end proof-assistant derivation.

\subsection{Root extraction before index or sign semantics}

For proof notation only put
\[
 \delta=Q-1,\quad D_y=QV^2-1,\quad z=Z/M,\quad
 \kappa=K_+/H,\quad \rho=e/D,\quad t=B/h,\quad W=1+hL.
\]
The identities
\begin{equation}
 P_L=\delta Z^2-M^2GD_y,\qquad
 z=n+\kappa-Wt-\rho
 \label{eq:compact-carrier}
\end{equation}
hold without divisibility assumptions. Positivity of the inputs gives
$q\ge32$, $X\ge q$, $Y\ge q^3$, $k\ge2$, $c\ge3$, $A>4$,
$\Delta\ge3$, $U\ge c^2$, $E>1$, $F_*>c+1$, $Q>1$, and $M>0$.
For every integer $b$, $V=F_*b-c\ne0$: at $b=0$ it is $-c$;
otherwise $|V|\ge F_*-c>1$. Therefore $D_y>0$.

Identically,
\[
 F_*=(1+2a)^2,\qquad Q=4\Delta^2U^2G.
\]
Thus $G$ is coprime to both $\delta$ and $D_y$. A zero of
\eqref{eq:compact-unscaled} gives $\delta z^2=GD_y$.
For each prime dividing $G$, its valuation in $G$ is even, because
the other two factors are units at that prime and a rational square
has even valuation. Hence $G=g^2$ with $g>0$ integral. Since
$F_0\equiv1\pmod E$, the positive coprime factors split:
\[
 E=\tau^2,\qquad F_0=f_0^2,\qquad g=\tau f_0.
\]
Then $y=|z|/g$ is a positive rational satisfying
\begin{equation}
 QV^2-(Q-1)y^2=1.
 \label{eq:compact-adjacent}
\end{equation}

For completeness, the special adjacent form forces $y$ integral.
Writing $y=r/m$ in lowest terms gives $m^2\mid Q-1$. If
$v=|V|>1$, then $v<y<Qv/(Q-1)$. The norm-preserving transformation
\[
 v'=(2Q-1)v-2(Q-1)y,\qquad
 y'=(2Q-1)y-2Qv
\]
has an integer first coordinate with $0<|v'|<v$ and preserves the
reduced denominator $m$ of its second coordinate. Reflecting signs and
descending reaches $v=1,y=1$, so $m=1$. Consequently
\[
 z=\varepsilon gy\in\Z,\qquad \varepsilon\in\{-1,1\}.
\]
Set $f=1+2a$ and $I=2Ug$. Pell doubling gives
\begin{equation}
 f^2-\Delta I^2=1,\qquad Q=(\Delta I)^2,\qquad c^2\mid I,
 \qquad V=f^2[c(T+\alpha)-R]-c.
 \label{eq:compact-strong}
\end{equation}
No residue involving a scale $\mu$ has been used.

\subsection{Main folding and transport signs}

The folding interface applies to \eqref{eq:compact-adjacent} and
\eqref{eq:compact-strong}, since $c\le I$ and $V\equiv-c\pmod f$.
It gives $c=\psi_A(p)$ for odd $p\ge1$. The first norm yields
$k=2\psi_{2V_0+1}(n_0)$ for $n_0\ge1$. If $n_0=1$, then
$k=2$, $\eta=\zeta=1$, and $c=2Y+1$ lies strictly between the first
two main ordinates $1$ and $2A$, impossible. Therefore
\begin{equation}
 k\ge8V_0+4,\qquad
 c>kY\ge k>8XY^2\ge8q^7.
 \label{eq:compact-large-c}
\end{equation}
This replaces the earlier scale-based $8q^5$ bound by a stronger one.
It proves $c>X$ and $K_+>0$.
The unconditional mask estimate
$R<q^4+(\alpha+2)q^2$, together with
\eqref{eq:compact-large-c}, gives $b>c$ before any sign is assigned
to $B$. Thus $V>0$, $\kappa>0$, and
\begin{equation}
 b-e=c(T+\alpha)-k+1>0,\qquad \rho<b/D.
 \label{eq:compact-rho}
\end{equation}

For $L=V^2$, the identity
\[
 a^2-g^2=E\bigl[\Delta U^2E(\Delta U^2-1)-1\bigr]>0
\]
gives $g<a$, whence $V>(F_*-1)b>ab>3g$. Since
$y<3V/2$ by \eqref{eq:compact-adjacent}, $gy<L/2$.

For $L=bV$, the same fifth-index argument as in package 133 gives
$b>4g$. In detail, $c\ge3$ makes $p\ge3$, so $c\ge4\Delta+3$ and
$V>4Q$. Put $S=\Delta I$, so $Q=S^2\ge4$. Classification of
\eqref{eq:compact-adjacent} gives $SV=\chi_S(j)$ with odd $j$.
The first two odd indices give $V=1,4Q-3$, both excluded. Hence
\[
 V\ge16Q^2-20Q+5>8Q^2,
 \qquad b>4\Delta^2(F_*-1)=4\Delta^3I^2>4g.
\]
It follows that $gy<3bV/8<L/2$ here as well.

If $B\le-1$, then $Wt\le-L-1/h$. Equations
\eqref{eq:compact-carrier} and \eqref{eq:compact-rho} imply
\[
 n=\varepsilon gy-\kappa+Wt+\rho
   <L/2-L+b/D<0,
\]
because $V>b>1$ and $D>2$. This contradicts $n>0$. Therefore
$B\ge0$ and $C>0$. The outer mask bounds now give
$3q<R<q^4-2$. The rank interface has all its premises:
$c^2\mid I$, $c>2R$, $c>2p$, $V\equiv-c\pmod f$, and
$V\equiv-R\pmod c$. It gives $p=R$; the unchanged mask parity
then yields even $q$ and odd $J$.

\subsection{The new staged index lemma}

\begin{lemma}[Partial divisibility recovers the full denominator]
At every positive zero of \eqref{eq:compact-unscaled},
\[
 n_0=(R+1)/2\qquad\text{and}\qquad Y^2\mid e.
\]
Neither $\gcd(Y,q-1)=1$ nor integrality of $e/Y^2$ is needed as a premise.
\end{lemma}

\begin{proof}
The integral carrier in \eqref{eq:compact-carrier}, after clearing
$hH$, implies $Y^2\mid hHe$. Since $Y=q^3s$, one has $q^6\mid Y^2$.
Moreover $h=q-1$ and $H=4Y(X+1)-1\equiv-1\pmod q$, so
$\gcd(q^6,hH)=1$. Consequently $q^6\mid e$.

Let $P_0=2XY^2+1$. The base-one Pell recurrence gives
\[
 \psi_{P_0}(n_0)\equiv n_0\pmod{P_0-1},\qquad
 k\equiv2n_0\pmod{4XY^2}.
\]
It follows that $q^6\mid2n_0-R-1$.
But $P_0-A=Y(2XY-X-1)>0$. If $n_0\ge R$, monotonicity in the
base and index gives $k/2=\psi_{P_0}(n_0)\ge\psi_A(R)=c$,
contradicting $c=kY+\eta>k$. Thus $1\le n_0<R$, and
\[
 |2n_0-R-1|<R<q^4<q^6.
\]
The divisible integer must be zero. Returning to the \emph{full}
first-Pell congruence now gives $4XY^2\mid k-R-1=e$, and in
particular $Y^2\mid e$.
\end{proof}

The modulus $q^6$ is used only in this proof. It introduces no new
polynomial register. The whole denominator is recovered after the
index, which is the step that removes the artificial scale.

\subsection{Projection, source recovery, and positive completeness}

Write $R=2r+1$, with $r\ge48$. The retained zero-offset ratio lemma
holds for every $\lambda\ge1$ with
$A_\lambda=\lambda Y(X+1)+1$ and
$P_\lambda=2X\lambda^2Y^2+1$. At $\lambda=1$ it gives
\[
 \xi<\frac{\psi_A(2r+1)}{\psi_{P_0}(r+1)}
 <\xi\left(1+\frac1{Y(X+1)}\right)^{2r},\qquad
 \xi=\frac{(X+1)^{2r}}{X^r}.
\]
The lemma explicitly includes this value of $\lambda$. Its lower
comparison follows from
\[
 (2A_\lambda-1)^2X-(X+1)^2(2P_\lambda)
 =4\lambda Y(X+1)X+X-2(X+1)^2>0;
\]
the upper comparison uses
$A_\lambda\le\lambda(Y(X+1)+1)$.
Together with $Y<c/k<Y+1$, it yields
$Y>X^r/3$ and $X^{r+1}<3Y(X+1)<H$.

Now $\rho$ is integral, so the carrier gives $H\mid hK_+$.
The identity
\[
 K_+=HXc-2K_{\mathrm{quad}},\qquad
 K_{\mathrm{quad}}=X[X+2(A-2)c]-1
\]
allows cancellation of two through odd $H$. The universal projection
therefore implies
\[
 H\mid h(X-2^R)(2^RX+1).
\]
The retained whole-product bound applies with the same $X\ge q\ge32$
and $r\ge48$; its old multiplier $q^3h$ was larger. Thus the entire
product has absolute value below $X^{r+1}<H$, and $X=2^R$.
No cancellation of a possibly nonunit factor modulo a composite $H$
is assumed. It follows that $\kappa$ is integral. The remaining carrier
gives $h\mid WB$, hence $h\mid B$ because $W\equiv1\pmod h$.
The quotient $t\ge0$ is integral and $q^3\mid Y$ is literal.
Fixed-minus parity and exact half-binomial rounding now reconstruct
the earlier source with its \emph{fresh} base $Y(X+1)+2$.
The pinned native calibration yields a finite ZFC refutation.

Conversely, a refutation supplies canonical outer data with
$X=2^R$, $R\equiv3\pmod4$, $r\ge48$, $q^3\mid Y$, and
$t=B/h\ge0$, together with the half-binomial identity. Set $s=Y/q^3>0$.
At the new unscaled bases take
\[
 c=\psi_A(R),\quad k=2\psi_{P_0}(r+1),\quad
 \tau=\chi_{P_0}(r+1),\quad
 \eta=c-kY,\quad\zeta=k-\eta.
\]
The uniform ratio bound and the same rounding estimates prove
$0<\eta<k$, so both new coordinates are positive.
The base-one recurrence gives $e$ a positive multiple of $4XY^2$,
so $\rho\ge1$; projection gives a positive integral $\kappa$.

Choose a seed Pell pair $(f_0,I_0)$ at a positive multiple of the
finite matrix period modulo $c^2\tau$, and set
$i=I_0/(c^2\tau)>0$. Then $U=I_0/\tau$, $F_0=f_0^2$, and
$G=(\tau f_0)^2$. Double this pair:
\[
 f=2f_0^2-1,\qquad I=2f_0I_0=2Ug.
\]
The retained auxiliary-completeness theorem accepts this strong
pair without a restriction on the parity of its index. It supplies
$T'=T+\alpha>\alpha$ and positive $V,y$ satisfying the stated auxiliary
equation, so $T>0$. Finally choose
\[
 n=gy+Wt+\rho-\kappa.
\]
Every term is integral. Since
$f_0>(A-1)I_0>Xc>\kappa$, one has $g>\kappa$; hence $n>0$.
The resulting carrier is exactly $z=gy$, and the polynomial vanishes.
This constructs all ten positive inputs freshly and proves both
existential directions.

\subsection{Exact cost, degree, and checks}

The actual circuits have the following operation counts:
\begin{center}
\begin{tabular}{lrrrrr}
\toprule
Weight&Multiplications&Additions&Subtractions&Total&Degree\\
\midrule
$V^2$&46&18&14&78&630\\
$bV$&47&18&14&79&468\\
\bottomrule
\end{tabular}
\end{center}
The relevant global degrees are
\[
 \begin{array}{c|rrrrrrrrr}
 \text{register}&A&\Delta&E&U&a&F_*&Q&G&V\\\hline
 \deg&6&12&22&11&56&112&124&78&118
 \end{array}
\]
with $\deg b=6$, $\deg M=15$, and $\deg B=2$.
For $L=V^2$, $\deg Z=253$ and $QZ^2$ uniquely supplies degree 630.
For $L=bV$, $\deg Z=141$ and $-QM^2GV^2$ uniquely supplies degree 468.

Let $q_0=\beta J$, $k_0=\eta+\zeta$, and $t_0=T+\alpha$.
The entire leading homogeneous components are
\begin{align*}
 [P_{V^2}]_{630}
 &=2^{14}q_0^{382}w^{46}s^{112}i^{20}k_0^{64}t_0^4
       (q_0-F-\alpha)^2,\\
 [P_{bV}]_{468}
 &=-2^{10}q_0^{288}w^{34}s^{84}i^{14}k_0^{46}t_0^2.
\end{align*}
They are nonzero at the actual fixed positive $\beta$.
In particular the displayed formulas contain coefficients
$2^{14}\beta^{382}$ and $-2^{10}\beta^{288}$ on the respective
monomials
\[
 J^{382}w^{46}s^{112}i^{20}\eta^{64}T^4F^2,
 \qquad J^{288}w^{34}s^{84}i^{14}\eta^{46}T^2.
\]
These also illustrate the cost of the compact representation: the
expanded coefficients can be much larger than the stored ports.

The independent standard-library checker proves 28 exact polynomial
cut identities per circuit, covering all 78 or 79 actual registers.
It verifies topology, backward liveness, every fixed coefficient byte,
the complete leading forms (1,950 and 141 terms), and eight actual
algebraic corruptions per circuit. Normal and optimized Python checks
and byte-identical circuit rebuilds are supplied separately.
These finite checks verify the artifacts and their algebra; the
unbounded zero-existence equivalence is the ordinary proof above and
its stated inherited interfaces.
